\section{Research Agenda}
\label{sec:agenda}

The theory generates seven immediate research programs, each derived from a specific proposition and each producing falsifiable predictions.

\textbf{1. Cross-domain morphospace replication} (from Proposition 1). The theory predicts that any class of designed artifacts will exhibit the same morphospace structure: non-uniform occupation, three-region partition, and cumulative expansion. Immediate targets include building facades (architectural morphospace), poster typography (graphic design morphospace), and software interface layouts (interaction design morphospace). A single replication that confirms the three-region structure in a non-furniture domain would substantially strengthen the theory's generality claim.

\textbf{2. Style causality experiment} (from Proposition 6). No existing experiment isolates style exposure as an independent variable on form variation. The theory predicts that style is a pattern memory, not a causal force: exposure to a style constrains the designer's cognitive cone but does not generate selection pressure. The falsification experiment: give two groups of designers identical functional briefs and material constraints but different style exposure (one group sees Art Nouveau chairs, the other sees Bauhaus chairs). If the resulting forms differ \textit{only} in the dimensions corresponding to the style exemplars; and not in the dimensions corresponding to the shared pressures; the theory holds. If style exposure shifts form along functional dimensions, the theory requires revision.

\textbf{3. Care dimensionality metric} (from Proposition 5). The theory defines care as the number of morphospace axes the agent actively represents, but does not yet provide an operational measurement protocol. A think-aloud protocol analysis of designers during design sessions could yield such a metric: count the distinct constraint dimensions verbalized during the session. The prediction: designers who verbalize more dimensions produce forms that score higher on multi-criteria robustness evaluations.

\textbf{4. Agentic design tool prototype} (from Proposition 6). The theory implies a specific architecture for AI-assisted design tools: the human sets the cognitive cone parameters (which dimensions to attend to, which constraints to enforce), and the AI executes shape grammar rules within that cone. A prototype implementation; where the designer specifies care dimensions via an explicit interface and the generative system explores within those dimensions; would test whether cone-width correlates with design quality as measured by expert evaluation.

\textbf{5. Non-European independent convergence} (from Proposition 1). The theory predicts that landscape attractors are real, not culturally contingent. If Chinese, Japanese, and West African chair traditions converge on the same morphospace attractors as European traditions; despite independent development; the attractor structure is confirmed as a physical-ergonomic reality. If convergence is absent, the theory must accommodate culturally specific landscape features.

\textbf{6. Material latency quantification} (from Proposition 2). The theory predicts a systematic lag between material availability and full formal exploitation. This is testable across multiple material transitions: wood to steel (lag: approximately 400 years from structural steel to tube chair), steel to plastic (lag: approximately 50 years from Bakelite to monobloc), plastic to composites, composites to biomaterials. If the latency pattern is consistent across transitions; long initial lag, rapid exploitation once the operation is discovered, then plateau; the theory's latency mechanism is confirmed.

\textbf{7. Multi-species design formalism} (from Proposition 5). The theory's agent definition is substrate-independent, but the worked examples (Section~\ref{sec:examples}) reveal a practical challenge: multi-species design requires coordinating agents whose cognitive cones operate on incommensurable timescales. A formal extension of the agent definition to accommodate incommensurable cones; with empirical testing against actual multi-species design projects such as constructed wetlands, permaculture systems, or urban rewilding initiatives; would test the theory's scalability beyond human-centered design.

Each of these programs is independently publishable, empirically tractable, and directly derived from the theory's propositions. Together, they constitute a decade-long research agenda that would either confirm the theory's generality or identify the specific conditions under which it requires revision.
